\chapter{Actas de reunión}\label{anexo:actas}

\section*{Acta de Reunión 1}
\addcontentsline{toc}{section}{Acta de Reunión 1}

\subsection*{Información General}
\begin{itemize}
    \item \textbf{Fecha:} 14 de julio de 2025
    \item \textbf{Hora:} 11:00 a.m.
    \item \textbf{Duración:} 30 minutos
    \item \textbf{Lugar:} Escuela Técnica Superior de Ingeniería Informática
\end{itemize}

\textbf{Participantes:}
\begin{itemize}
    \item José Antonio Parejo Maestre (Tutor)
    \item Cynthia Castaño Juan (Estudiante)
\end{itemize}

\subsection*{Objetivo de la reunión}
Presentar y proponer al tutor el tema del Trabajo de Fin de Grado (TFG) y discutir las dimensiones principales del proyecto (alcance funcional, complejidad técnica y viabilidad en el tiempo).

\subsection*{Temas tratados}

\subsubsection*{1. Propuesta del tema del TFG}
La estudiante propone como tema del TFG el desarrollo de una aplicación para la gestión y consulta de eventos culturales y de ocio, centrada inicialmente en una ciudad concreta.

El tutor valora positivamente la propuesta y confirma que encaja con los objetivos formativos del grado y con su área de tutoría.

\subsubsection*{2. Definición de las dimensiones del proyecto}
Se discuten las dimensiones principales del TFG:

\begin{itemize}
    \item \textbf{Dimensión funcional:} alcance inicial de funcionalidades (búsqueda y visualización de eventos, fichas de detalle, posibles opciones futuras como mapa, reseñas o recomendaciones).
    \item \textbf{Dimensión técnica:} tecnologías potenciales a utilizar (aplicación móvil o web, arquitectura general, posibles servicios externos).
    \item \textbf{Dimensión temporal:} viabilidad del proyecto dentro del calendario del TFG, estimando el esfuerzo necesario y la necesidad de acotar el alcance para garantizar su finalización.
\end{itemize}

\subsubsection*{3. Viabilidad y acotación inicial}
Se comenta la importancia de mantener un alcance realista para asegurar que el proyecto pueda completarse con la calidad requerida dentro del plazo establecido.

Se acuerda que en reuniones posteriores se detallará mejor el backlog y la planificación, una vez formalizado el tema.

\subsection*{Acuerdos y decisiones}
\begin{enumerate}
    \item El tutor acepta de forma preliminar el tema propuesto para el TFG, pendiente de la formalización administrativa correspondiente.
    \item Se acuerda acotar el proyecto a un conjunto de funcionalidades básicas para garantizar su viabilidad (definición detallada en siguientes reuniones).
    \item La estudiante elaborará una primera propuesta de descripción del TFG, incluyendo objetivos, alcance y dimensiones del proyecto, para compartir con el tutor.
\end{enumerate}

\subsection*{Tareas asignadas}

\textbf{Estudiante (Cynthia Castaño Juan)}
\begin{itemize}
    \item Redactar un breve documento de propuesta del TFG (título provisional, objetivos, justificación y alcance inicial).
    \item Enviar la propuesta al tutor para revisión antes de la siguiente reunión.
\end{itemize}

\textbf{Tutor (José Antonio Parejo Maestre)}
\begin{itemize}
    \item Revisar la propuesta de tema y alcance del TFG.
    \item Orientar sobre posibles ajustes en las dimensiones (funcional, técnica y temporal) en la próxima reunión.
\end{itemize}

\subsection*{Próxima reunión}
\begin{itemize}
    \item \textbf{Objetivo:} Revisar y, en su caso, validar la propuesta escrita del TFG (tema, objetivos y alcance) y definir los siguientes pasos de planificación.
    \item \textbf{Fecha:} Por determinar.
\end{itemize}

\textbf{Elaborado por:} Cynthia Castaño Juan \hfill \textbf{Fecha:} 3 de enero de 2026

\newpage

\section*{Acta de Reunión 2}
\addcontentsline{toc}{section}{Acta de Reunión 2}

\subsection*{Información General}
\begin{itemize}
    \item \textbf{Fecha:} 25 de septiembre de 2025
    \item \textbf{Hora:} 10:00 a.m.
    \item \textbf{Duración:} 30 minutos
    \item \textbf{Lugar:} Reunión online (videollamada)
\end{itemize}

\textbf{Participantes:}
\begin{itemize}
    \item José Antonio Parejo Maestre (Tutor)
    \item Cynthia Castaño Juan (Estudiante)
\end{itemize}

\subsection*{Objetivo de la reunión}
Revisar y discutir el documento de Anteproyecto de TFG ``Aplicación móvil de eventos emergentes en Sevilla, `Sevillaneando''', contrastando la descripción, objetivos y tecnologías iniciales, y acordar una ampliación y concreción del alcance funcional del proyecto.

\subsection*{Temas tratados}

\subsubsection*{1. Revisión del anteproyecto}
Se revisa la descripción del proyecto, centrada en una aplicación móvil para promover eventos culturales y de ocio emergentes en Sevilla, incluyendo también eventos convencionales.

El tutor considera adecuada la motivación y confirma que la propuesta encaja con el ámbito de Ingeniería del Software y con la línea de proyectos de aplicaciones móviles.

\subsubsection*{2. Alcance funcional y funcionalidades iniciales}
Se comentan los objetivos del anteproyecto, que incluyen: listado de eventos con información detallada, posibilidad de que los usuarios añadan eventos sujetos a moderación, integración de mapa, subida de imágenes en vivo, reseñas y valoraciones, recomendaciones personalizadas y notificaciones push.

A partir de estos objetivos, se acuerda estructurar el alcance en bloques funcionales (búsqueda/listado, detalle de evento, propuesta y moderación, mapa, interacción social, recomendaciones y notificaciones), para facilitar su futura transformación en historias de usuario.

\subsubsection*{3. Matización del alcance}
Durante la reunión se decide:

\begin{itemize}
    \item Confirmar el uso de un sistema de moderación de contenido para fotos en vivo, basado en un modelo propio con NSFW.js/TensorFlow.js integrado en el backend, considerándolas funcionalidades avanzadas pero incluidas en el TFG.
\end{itemize}

\subsubsection*{4. Validación de tecnologías iniciales}
Se revisan las decisiones tecnológicas recogidas en el anteproyecto: React Native para el frontend móvil, backend con Node.js y NestJS, PostgreSQL con PostGIS para datos geoespaciales, Firebase Auth y Firebase Storage, y Google Maps SDK para mapas.

El tutor valida estas decisiones como técnicamente coherentes y adecuadas para el tipo de aplicación propuesta, recomendando documentar claramente las razones de elección en la memoria del TFG.

\subsection*{Acuerdos y decisiones}
\begin{enumerate}
    \item El documento de Anteproyecto queda aceptado como base oficial para el desarrollo del TFG ``Sevillaneando'', con la descripción, objetivos y tecnologías recogidas en él.
    \item Se amplía y concreta el alcance funcional, fijando como parte del proyecto:
    \begin{itemize}
        \item Listado y detalle de eventos con información completa.
        \item Alta de eventos por usuarios con flujo de moderación.
        \item Integración de mapa con búsqueda por proximidad.
        \item Subida de imágenes en vivo con moderación automática.
        \item Sistema de reseñas y valoraciones.
        \item Módulo de recomendaciones personalizadas y notificaciones push.
    \end{itemize}
    \item Se acuerda utilizar este alcance como base para derivar el backlog inicial de épicas e historias de usuario.
\end{enumerate}

\subsection*{Tareas asignadas}

\textbf{Estudiante (Cynthia Castaño Juan)}
\begin{itemize}
    \item Actualizar el anteproyecto, si es necesario, incorporando claramente el listado de funcionalidades acordadas.
    \item Elaborar un primer backlog de historias de usuario alineado con los objetivos y tecnologías del anteproyecto.
\end{itemize}

\textbf{Tutor (José Antonio Parejo Maestre)}
\begin{itemize}
    \item Revisar el backlog inicial y proponer ajustes de alcance o prioridades en la siguiente reunión.
\end{itemize}

\subsection*{Próxima reunión}
\begin{itemize}
    \item \textbf{Objetivo:} Revisar el backlog inicial de historias de usuario, validar la coherencia entre alcance, funcionalidades y tecnologías, y esbozar una primera planificación por sprints.
    \item \textbf{Fecha:} Por determinar (a concretar por correo tras la entrega del backlog).
\end{itemize}

\textbf{Elaborado por:} Cynthia Castaño Juan \hfill \textbf{Fecha:} 3 de enero de 2026

\newpage

\section*{Acta de Reunión 3}
\addcontentsline{toc}{section}{Acta de Reunión 3}

\subsection*{Información General}
\begin{itemize}
    \item \textbf{Fecha:} 10 de octubre de 2025
    \item \textbf{Hora:} 9:58 a.m.
    \item \textbf{Duración:} 30 minutos 13 segundos
\end{itemize}

\textbf{Participantes:}
\begin{itemize}
    \item José Antonio Parejo Maestre (Tutor)
    \item Cynthia Castaño Juan (Estudiante)
\end{itemize}

\subsection*{Objetivo de la reunión}
Revisión del backlog del proyecto, validación del modelo conceptual de datos, definición de la planificación inicial y establecimiento de objetivos para el siguiente sprint.

\subsection*{Temas tratados}

\subsubsection*{1. Aprobación del Product Backlog}
Se revisó y aprobó el backlog inicial presentado. Se acordó que el backlog es un documento vivo que se irá modificando y perfilando durante todo el proceso de desarrollo. Se considera suficiente contar con una primera versión de historias básicas bien revisadas para arrancar la fase de desarrollo.

\subsubsection*{2. Reorganización de la planificación}
Se acordó reorganizar la estructura de fases:

\begin{itemize}
    \item \textbf{Fase 1:} Incluye tareas 1.1, 1.2, 1.3 y la creación del modelo conceptual de datos.
    \item \textbf{Fase 2:} El diseño de arquitectura (tarea 1.4) pasa a la fase 2 (2.01).
\end{itemize}

Se confirmó el stack tecnológico: React Native para el frontend y backend con Node.js (Next.js).

\subsubsection*{3. Estimación de tareas}
Se definió utilizar la técnica de sizing con tallas de camiseta (XS, S, M, L, XL) para estimar el esfuerzo de las historias de usuario y tareas relevantes. Este método utiliza medidas relativas en lugar de horas fijas.

\subsubsection*{4. Revisión del modelo conceptual (UML)}
Se revisó el diagrama de clases del modelo conceptual con las siguientes entidades principales:

\textbf{Entidad Evento:}
\begin{itemize}
    \item Atributos: ID, título, descripción, fecha/hora inicio, fecha/hora fin, ubicación (coordenadas), precio, categorías, estado (pendiente/aprobado/rechazado), creador, imágenes, reseñas.
    \item Se aclaró que fecha inicio/fin son instantes en el tiempo, no solo fechas.
    \item Las imágenes son referencias URL almacenadas en el servidor.
\end{itemize}

\textbf{Entidad Mapa:}
\begin{itemize}
    \item Se acordó añadir campos adicionales: nivel de zoom, centro (latitud/longitud) o esquinas (punto inicio/fin).
    \item Esto permite controlar la visualización de eventos relacionados en una zona específica.
\end{itemize}

\textbf{Nueva Entidad: Ruta}
\begin{itemize}
    \item Incluye: conjunto de coordenadas, secuencia de eventos, temporización.
    \item Permite planificar rutas por múltiples eventos en un día o periodo concreto.
\end{itemize}

\textbf{Otras entidades:}
\begin{itemize}
    \item Categorías, Imágenes, Reseñas.
    \item Usuario con tipo de usuario (incluyendo Moderador).
\end{itemize}

\textbf{Entidad Recomendación:}
\begin{itemize}
    \item Un conjunto de eventos recomendados automáticamente a un usuario específico.
    \item Se acordó añadir un flag para indicar si el usuario ya ha visto la recomendación.
    \item Requiere definir un algoritmo de generación basado en criterio de similaridad (inicialmente por ubicación).
\end{itemize}

\textbf{Entidad Notificación:}
\begin{itemize}
    \item Debe relacionarse con recomendaciones y eventos guardados.
    \item Requiere mecanismos de configuración para usuarios.
    \item Debe incluir información del evento específico.
\end{itemize}

\subsubsection*{5. Metodología de desarrollo}

\textbf{Control de versiones:}
\begin{itemize}
    \item Git con política de ramas: Feature Branch (3 niveles: main/develop/feature).
    \item Mensajes de commit siguiendo Conventional Commits.
\end{itemize}

\textbf{Gestión del proyecto:}
\begin{itemize}
    \item Se utilizará una adaptación de Scrum para TFG:
    \begin{itemize}
        \item No se realizan daily standups.
        \item Las reuniones de review y planning se fusionan en una sola sesión.
        \item Se mantiene estructura de sprints con objetivos definidos.
    \end{itemize}
\end{itemize}

Se acordó completar la sección de metodología en el documento correspondiente.

\subsection*{Decisiones tomadas}
\begin{enumerate}
    \item El backlog inicial queda aprobado y se actualizará durante el desarrollo.
    \item Se incorpora la entidad ``Ruta'' al modelo conceptual.
    \item Se añaden campos de visualización (zoom, centro) a la entidad ``Mapa''.
    \item Se utilizará sizing con tallas de camiseta para estimación.
\end{enumerate}

\subsection*{Tareas asignadas}

\textbf{Estudiante (Cynthia):}
\begin{itemize}
    \item Desarrollar Producto Mínimo Viable (MVP) con sistema de login/logout y visualización de mapa en pantalla.
    \item Actualizar el modelo conceptual incorporando las modificaciones discutidas.
    \item Completar la sección de metodología en el documento de gestión.
    \item Crear repositorio del proyecto con la política de ramas acordada.
    \item Aplicar estimaciones con sizing a las historias y tareas.
    \item Conceder acceso al tutor del sistema de seguimiento de horas del proyecto.
\end{itemize}

\textbf{Tutor (José Antonio):}
\begin{itemize}
    \item Realizar adjudicación oficial del TFG antes de mayo de 2026.
\end{itemize}

\subsection*{Próxima reunión}

Objetivos para el siguiente sprint:

\begin{enumerate}
    \item Presentar MVP funcional con login/logout y visualización de mapa.
    \item Mostrar modelo conceptual actualizado.
    \item Backlog actualizado con historias de notificaciones y recomendaciones.
    \item Documento de metodología completado.
    \item Repositorio creado y configurado.
\end{enumerate}

\textbf{Fecha:} Por determinar

\textbf{Elaborado por:} Cynthia Castaño Juan \hfill \textbf{Fecha:} 2 de enero de 2026

\newpage

\section*{Acta de Reunión 4}
\addcontentsline{toc}{section}{Acta de Reunión 4}

\subsection*{Información General}
\begin{itemize}
    \item \textbf{Fecha:} 9 de enero de 2026
    \item \textbf{Hora:} 13:25
    \item \textbf{Duración:} 32 minutos 41 segundos
\end{itemize}

\textbf{Participantes:}
\begin{itemize}
    \item José Antonio Parejo Maestre (Tutor)
    \item Cynthia Castaño Juan (Estudiante)
\end{itemize}

\subsection*{Objetivo de la reunión}
Revisión del estado del proyecto, validación del modelo conceptual, definición de prioridades del backlog y planificación del siguiente sprint.

\subsection*{Temas tratados}

\subsubsection*{1. Revisión del modelo conceptual}
Se revisó el modelo conceptual actualizado, especialmente las entidades relacionadas con eventos, imágenes, usuarios y reseñas. Se aclaró que:

\begin{itemize}
    \item Una imagen pertenece a un único usuario (el que la sube), aunque pueda estar asociada a múltiples eventos.
    \item Un evento puede tener múltiples imágenes.
    \item Las relaciones entre eventos, reseñas y usuarios quedan correctamente definidas.
    \item No es necesario vincular rutas directamente a eventos, ya que la recomendación por proximidad se obtiene mediante la estructura geográfica.
\end{itemize}

Se confirmó que el sistema utilizará coordenadas geográficas (latitud y longitud) para representar ubicaciones.

\subsubsection*{2. Sistema de mapas}
Se acordó utilizar OpenStreetMap en lugar de Google Maps debido a sus menores restricciones y mayor facilidad de integración.

Se almacenará la ubicación mediante un GeoPoint (latitud, longitud) para permitir cálculos de distancia y filtrado geográfico.

Esta decisión debe documentarse como una decisión de diseño.

\subsubsection*{3. Estado actual del proyecto}
Se confirmó que actualmente están implementadas las siguientes funcionalidades:

\begin{itemize}
    \item Registro y gestión de usuarios.
    \item Tres tipos de usuario: Usuario, Moderador y Administrador.
    \item Edición de perfil, cambio de contraseña y eliminación de cuenta.
    \item Creación de eventos por usuarios.
    \item Moderación de eventos por moderadores.
    \item Visualización de eventos en mapa.
    \item Gestión de categorías (por organizadores).
    \item Notificaciones cuando un evento es aprobado o rechazado.
\end{itemize}

Se indicó que la edición de eventos por moderadores aún está pendiente.

\subsubsection*{4. Backlog y priorización}
Se acordó sustituir la priorización Alta / Media / Baja por el modelo MoSCoW (Must, Should, Could, Won't), ya que proporciona mayor claridad sobre qué funcionalidades son esenciales y cuáles son opcionales.

\subsubsection*{5. Metodología y documentación}
Se confirmó que ya están definidos: estimación por tallas de camiseta, políticas de ramas, políticas de versionado, política de commits y planificación tipo Gantt.

Se detectó como pendiente añadir una guía de estilo de codificación, incluyendo reglas de formato, convenciones y configuración del entorno para autoformateo.

\subsubsection*{6. Funcionalidades opcionales propuestas}

\textbf{Generación automática de carteles:} uso de una API de IA para generar carteles a partir de la descripción de los eventos, que se asociarán automáticamente al evento.

\textbf{Gestión avanzada de imágenes:}
\begin{itemize}
    \item Posibilidad de subir varias imágenes por evento.
    \item Límite de almacenamiento por usuario (ej. 15 MB).
    \item Proceso batch que calcule el uso y bloquee subidas al superar el límite.
\end{itemize}

\textbf{Sistema de chat:} implementación de chat entre usuarios mediante WebSockets, con una primera fase en tiempo real y una segunda con persistencia en base de datos.

\subsubsection*{7. Búsqueda y visualización de eventos}
Se acordó:

\begin{itemize}
    \item En la vista de mapa, mostrar eventos cercanos usando filtros por distancia geográfica.
    \item En la vista de búsqueda, permitir búsquedas más precisas por nombre, categoría o fecha.
\end{itemize}

\subsubsection*{8. Planificación del proyecto}
Se acordó que el siguiente sprint tendrá una duración de 4 semanas.

Las prioridades serán: mejora de la búsqueda y visualización geográfica, gestión de imágenes de eventos, edición de eventos por moderadores, gestión de categorías e inicio del sistema de chat.

Las funcionalidades de recomendaciones y scraping se dejan para sprints posteriores.

\subsection*{Decisiones tomadas}
\begin{itemize}
    \item Se adopta OpenStreetMap como sistema de mapas.
    \item Se utilizará GeoPoint (latitud/longitud) para ubicaciones.
    \item Se cambia la priorización del backlog a MoSCoW.
    \item Se incluyen como opcionales la generación de carteles, el chat y la gestión avanzada de imágenes.
    \item El sprint tendrá una duración de 4 semanas.
\end{itemize}

\subsection*{Tareas asignadas}

\textbf{Estudiante (Cynthia):}
\begin{itemize}
    \item Actualizar el backlog usando MoSCoW.
    \item Añadir la guía de estilo de codificación a la metodología.
    \item Implementar mejoras en búsqueda y visualización de eventos.
    \item Implementar edición de eventos por moderadores.
    \item Avanzar en la gestión de imágenes.
    \item Iniciar el sistema de chat.
\end{itemize}

\textbf{Tutor (José Antonio):}
\begin{itemize}
    \item Revisar los avances en la siguiente reunión.
    \item Validar las decisiones de diseño y priorización.
\end{itemize}

\subsection*{Próxima reunión}
Objetivos para el siguiente sprint:

\begin{itemize}
    \item Sistema de eventos con búsqueda geográfica funcional.
    \item Gestión de imágenes por evento.
    \item Edición de eventos por moderadores.
    \item Backlog priorizado con MoSCoW.
    \item Metodología actualizada con guía de estilo.
    \item Restricciones en el backend.
\end{itemize}

\textbf{Fecha:} Por determinar

\textbf{Elaborado por:} Cynthia Castaño Juan \hfill \textbf{Fecha:} 9 de enero de 2026

\newpage

\section*{Acta de Reunión 5}
\addcontentsline{toc}{section}{Acta de Reunión 5}

\subsection*{Información General}
\begin{itemize}
    \item \textbf{Fecha:} 17 de febrero de 2026
    \item \textbf{Hora:} 08:14
    \item \textbf{Duración:} 31 minutos 43 segundos
\end{itemize}

\textbf{Participantes:}
\begin{itemize}
    \item José Antonio Parejo Maestre (Tutor)
    \item Cynthia Castaño Juan (Estudiante)
\end{itemize}

\subsection*{Objetivo de la reunión}
Revisión del progreso del sprint anterior (mejoras de búsqueda, gestión de imágenes, edición por moderadores y chat), validación de la metodología de documentación y planificación del nuevo sprint enfocado en scraping.

\subsection*{Temas tratados}

\subsubsection*{1. Estado actual del proyecto y sprint finalizado}
Se validó la implementación de las funcionalidades acordadas en el sprint de cuatro semanas:

\begin{itemize}
    \item \textbf{Gestión de Categorías:} Finalizada. Solo los administradores pueden crearlas, aunque los usuarios pueden asociarlas a sus eventos.
    \item \textbf{Edición de Eventos por Moderadores:} Finalizada. Los moderadores pueden editar campos y aceptar o rechazar eventos.
    \item \textbf{Gestión de Imágenes:} Implementada la subida de una imagen de portada. Se discutió la migración a una tabla adicional para permitir múltiples imágenes por evento.
    \item \textbf{Búsqueda y Visualización Geográfica:} Implementado el mapa con filtros de proximidad codificados por colores (verde, amarillo, rojo) y enlace directo a Google Maps/Mapas, además de filtrados fuera del mapa interactivo por proximidad.
\end{itemize}

\subsubsection*{2. Sistema de chat}
Se presentó el sistema de chat en vivo utilizando WebSockets.

\begin{itemize}
    \item \textbf{Optimización de recursos:} El tutor advirtió sobre el consumo de memoria si los sockets permanecen abiertos indefinidamente.
    \item \textbf{Decisión técnica:} Se implementará una tarea programada (tipo cron) en el servidor para cerrar conexiones de eventos finalizados o chats inactivos tras un tiempo configurable.
\end{itemize}

\subsubsection*{3. Metodología y documentación}
\begin{itemize}
    \item \textbf{Guía de estilo:} Se confirmó la creación de una guía de estilo en formato Markdown y la inclusión de scripts de linting (ESLint/Prettier) para automatizar el cumplimiento de reglas en el backend y frontend.
\end{itemize}

\subsubsection*{4. Planificación del nuevo sprint}
Se definieron las prioridades para las próximas semanas:

\begin{itemize}
    \item \textbf{Scraping:} Extracción de datos de eventos de sitios clave (Turismo, Universidad, etc.).
    \item \textbf{Gestión de precios:} Se acordó añadir campos opcionales de ``precio máximo'' y ``precio mínimo'' para manejar eventos con costes variables o gratuitos.
    \item \textbf{Gestión de fotos:} Permitir múltiples imágenes por evento.
    \item \textbf{Gestión de chats:} Cerrar los sockets cuando se sale del chat y cuando se acaba el evento o tras un cierto tiempo.
\end{itemize}

\subsection*{Decisiones tomadas}
\begin{itemize}
    \item Se utilizará una tabla adicional en la base de datos para la colección de imágenes de eventos.
    \item Se implementará el cierre automático de websockets por inactividad o fin de evento para proteger el servidor.
    \item La memoria oficial se redactará en \LaTeX{} utilizando la plantilla de la Escuela.
    \item Se incorpora el uso de IA (Gemini) para facilitar el proceso de scraping y estructuración de datos.
\end{itemize}

\subsection*{Tareas asignadas}

\textbf{Estudiante (Cynthia):}
\begin{itemize}
    \item Migrar la lógica de imágenes a una colección/tabla dedicada.
    \item Implementar la lógica de cierre de chats en el backend.
    \item Implementar intervalo de precio.
    \item Preparar el primer borrador (índice y metodología) en Overleaf e invitar al tutor.
    \item Desarrollar los scrapers para fuentes de eventos institucionales.
    \item Actualizar el backlog y el diagrama UML.
\end{itemize}

\textbf{Tutor (José Antonio):}
\begin{itemize}
    \item Revisar los cambios y la estructura de la memoria en la siguiente sesión.
\end{itemize}

\subsection*{Próxima reunión}
\begin{itemize}
    \item \textbf{Objetivos:} Revisión de los datos extraídos mediante scraping, validación del primer borrador de la memoria.
    \item \textbf{Fecha:} Por determinar.
\end{itemize}

\textbf{Elaborado por:} Cynthia Castaño Juan \hfill \textbf{Fecha:} 17 de febrero de 2026

\newpage

\section*{Acta de Reunión 6}
\addcontentsline{toc}{section}{Acta de Reunión 6}

\subsection*{Información General}
\begin{itemize}
    \item \textbf{Fecha:} 27 de marzo de 2026
    \item \textbf{Hora:} 16:30
    \item \textbf{Duración:} 32 minutos
\end{itemize}

\textbf{Participantes:}
\begin{itemize}
    \item José Antonio Parejo Maestre (Tutor)
    \item Cynthia Castaño Juan (Estudiante)
\end{itemize}

\subsection*{Objetivo de la reunión}
Revisión del estado actual del proyecto tras el sprint en curso, definición del último sprint de desarrollo centrado en rutas de eventos y orientación sobre la fase de redacción de la memoria.

\subsection*{Temas tratados}

\subsubsection*{1. Estado actual del proyecto}
Se revisaron las funcionalidades desarrolladas durante el sprint actual y el cómputo total de horas invertidas:

\begin{itemize}
    \item Horas actuales: 160 horas.
    \item Se confirma que queda un único sprint de desarrollo.
\end{itemize}

\subsubsection*{2. Funcionalidades revisadas}
Se repasaron funcionalidades implementadas o en progreso:

\begin{itemize}
    \item Edición de eventos (públicos y privados).
    \item Calendario.
    \item Buscador.
    \item Notificaciones de eventos y recordatorios.
    \item Gestión de imágenes (posibilidad de múltiples fotos).
    \item Precio variable en eventos.
    \item Gestión de cierre de sockets.
\end{itemize}

Se discutieron mejoras opcionales: configuración del intervalo de notificaciones y mejora en gestión de imágenes (cropping).

\subsubsection*{3. Planificación del último sprint}
Se define como objetivo principal la implementación de rutas de eventos, incluyendo seguimiento de rutas dentro de la aplicación.

Como funcionalidades opcionales: mejoras en notificaciones y ajustes en funcionalidades existentes.

Se acuerda priorizar el objetivo principal y ajustar el alcance si el tiempo es limitado.

\subsubsection*{4. Organización del tiempo y entrega}
\begin{itemize}
    \item Se recomienda realizar un sprint aproximado de 3 semanas, ajustado por la feria (21--26 de abril).
    \item Próxima reunión propuesta: 17 de abril de 2026 (hora por confirmar).
\end{itemize}

\subsubsection*{5. Memoria y documentación}
Se enfatiza la importancia de comenzar la redacción de la memoria. Se identifican los tres manuales del TFG: manual de usuario, manual de desarrollo y manual de despliegue.

\subsubsection*{6. Testing y validación}
\begin{itemize}
    \item El testing de aceptación se realizará en paralelo con la redacción de la memoria.
    \item Se utilizarán estas pruebas para documentar la validación del sistema.
\end{itemize}

\subsubsection*{7. Dudas técnicas}

\textbf{Eventos pasados:} Se propone un toggle para mostrar/ocultar eventos pasados, o bien separarlos en una categoría/pestaña distinta.

\textbf{Despliegue de la aplicación:}
\begin{itemize}
    \item En iOS requiere cuenta de pago $\rightarrow$ se documentará en la memoria.
    \item Para pruebas: uso de emuladores Android y posibilidad de usar alternativas más ligeras que Android Studio.
\end{itemize}

\subsection*{Decisiones tomadas}
\begin{itemize}
    \item Priorizar la implementación de rutas de eventos en el último sprint.
    \item Iniciar la redacción intensiva de la memoria.
    \item Realizar el testing junto con la documentación.
    \item Mantener flexibilidad en funcionalidades opcionales según el tiempo disponible.
\end{itemize}

\subsection*{Tareas asignadas}

\textbf{Estudiante (Cynthia):}
\begin{itemize}
    \item Implementar el sistema de rutas de eventos.
    \item Avanzar en la redacción de la memoria.
    \item Elaborar los manuales (usuario, desarrollo y despliegue).
    \item Realizar pruebas de validación del sistema.
\end{itemize}

\textbf{Tutor (José Antonio):}
\begin{itemize}
    \item Revisar el progreso del sprint y la memoria en la próxima reunión.
    \item Proporcionar orientación sobre herramientas (emuladores, despliegue).
\end{itemize}

\subsection*{Próxima reunión}
\begin{itemize}
    \item \textbf{Fecha:} 17 de abril de 2026 (por confirmar hora).
    \item \textbf{Objetivos:} Revisión del sistema de rutas, evaluación del estado de la memoria y decisión sobre cierre de desarrollo o continuación.
\end{itemize}

\textbf{Elaborado por:} Cynthia Castaño Juan \hfill \textbf{Fecha:} 27 de marzo de 2026

\newpage

\section*{Acta de Reunión 7}
\addcontentsline{toc}{section}{Acta de Reunión 7}

\subsection*{Información General}
\begin{itemize}
    \item \textbf{Fecha:} 17 de abril de 2026
    \item \textbf{Hora:} 09:31
    \item \textbf{Duración:} 42 minutos
\end{itemize}

\textbf{Participantes:}
\begin{itemize}
    \item José Antonio Parejo Maestre (Tutor)
    \item Cynthia Castaño Juan (Estudiante)
\end{itemize}

\subsection*{Objetivo de la reunión}
Revisión del estado de la memoria del TFG, análisis de los manuales (usuario, desarrollo y despliegue), validación de diagramas de arquitectura y definición de tareas para el siguiente sprint centrado en pruebas, refactorización y cierre de la documentación.

\subsection*{Temas tratados}

\subsubsection*{1. Estado actual del proyecto}
Se revisa el progreso general del proyecto:

\begin{itemize}
    \item Horas estimadas actuales: entre 220 y 230 horas.
    \item Se confirma que el proyecto entra en fase final, con aproximadamente un mes hasta la entrega.
\end{itemize}

\subsubsection*{2. Revisión de manuales}
Se analizan los tres manuales del sistema (usuario, desarrollo y despliegue).

\textbf{Observaciones principales:}
\begin{itemize}
    \item El contenido es correcto y completo, pero demasiado esquemático.
    \item Se recomienda transformar el contenido en texto narrativo (más descriptivo y orientado a usuario).
    \item El manual de usuario debe centrarse en las necesidades del usuario y no en la estructura de la interfaz.
    \item Es necesario ampliar la introducción para dar más contexto.
\end{itemize}

\textbf{Mejoras propuestas:}
\begin{itemize}
    \item Añadir descripciones detalladas de pantallas, campos y validaciones.
    \item Incluir un grafo de navegabilidad de la aplicación.
    \item Incorporar una sección de errores frecuentes.
\end{itemize}

\subsubsection*{3. Estructura de la memoria}
\begin{itemize}
    \item Se recomienda estructurar los manuales como capítulos independientes dentro de la memoria.
    \item Se valida el uso de la EDT como base de planificación y se recomienda mantenerla simple (análisis, desarrollo y cierre).
    \item El presupuesto puede calcularse sobre estimación (300 horas) o sobre horas reales.
\end{itemize}

\subsubsection*{4. Diagramas de arquitectura}
Se revisan los diagramas realizados (vista de datos, vista de componentes/despliegue y diagramas de relaciones).

\textbf{Observaciones:}
\begin{itemize}
    \item El contenido es correcto, pero algunos diagramas no siguen la notación estándar UML.
    \item Se recomienda usar diagramas UML formales (clases, componentes, despliegue), mantener coherencia en la representación y mejorar la claridad visual.
\end{itemize}

\subsubsection*{5. Desarrollo y entrega}
\begin{itemize}
    \item Fecha de entrega: 20 de mayo de 2026.
    \item Se debe entregar: memoria y código fuente.
    \item El despliegue puede realizarse posteriormente de cara a la presentación.
    \item Es necesario tener una versión funcional completa del sistema.
\end{itemize}

\subsubsection*{6. Testing y validación}
Se establece como prioridad alta definir un plan de pruebas completo que incluya: historias de usuario a validar, datos de entrada, resultados esperados y resultados obtenidos. Se recomienda realizar capturas durante las pruebas para usarlas como evidencias en la memoria.

\subsubsection*{7. Próximo sprint}
Se define el objetivo del siguiente sprint (3 semanas):

\begin{itemize}
    \item Finalizar la refactorización del sistema.
    \item Diseñar y ejecutar el plan de pruebas.
    \item Completar la sección de validación.
    \item Mejorar y cerrar las secciones de la memoria.
\end{itemize}

\subsection*{Decisiones tomadas}
\begin{itemize}
    \item Priorizar la realización del plan de pruebas y validación del sistema.
    \item Reescribir los manuales en formato narrativo y no esquemático.
    \item Mejorar los diagramas siguiendo estándares UML.
    \item Considerar el desarrollo funcional prácticamente finalizado.
    \item Enfocar el trabajo restante en documentación y validación.
\end{itemize}

\subsection*{Tareas asignadas}

\textbf{Estudiante (Cynthia):}
\begin{itemize}
    \item Finalizar la refactorización del sistema.
    \item Elaborar el plan de pruebas.
    \item Ejecutar las pruebas y documentar resultados.
    \item Reescribir los manuales (usuario, desarrollo y despliegue).
    \item Mejorar diagramas de arquitectura.
    \item Completar y pulir la memoria.
\end{itemize}

\textbf{Tutor (José Antonio):}
\begin{itemize}
    \item Revisar avances de la memoria.
    \item Validar el plan de pruebas.
    \item Proporcionar feedback sobre la documentación.
\end{itemize}

\subsection*{Próxima reunión}
\begin{itemize}
    \item \textbf{Fecha:} 5 de mayo de 2026.
    \item \textbf{Objetivos:} Revisión del plan de pruebas y resultados, validación de la refactorización y evaluación del estado final de la memoria.
\end{itemize}

\textbf{Elaborado por:} Cynthia Castaño Juan \hfill \textbf{Fecha:} 17 de abril de 2026

\newpage

\section*{Acta de Reunión 8}
\addcontentsline{toc}{section}{Acta de Reunión 8}

\subsection*{Información General}
\begin{itemize}
    \item \textbf{Fecha:} 5 de mayo de 2026
    \item \textbf{Hora:} 10:31
    \item \textbf{Duración:} 35 minutos
\end{itemize}

\textbf{Participantes:}
\begin{itemize}
    \item José Antonio Parejo Maestre (Tutor)
    \item Cynthia Castaño Juan (Estudiante)
\end{itemize}

\subsection*{Objetivo de la reunión}
Revisión detallada del borrador completo de la memoria del TFG, identificación de mejoras estructurales y de contenido, validación de aspectos metodológicos y definición de ajustes finales antes de la entrega.

\subsection*{Temas tratados}

\subsubsection*{1. Estado actual del proyecto}
\begin{itemize}
    \item La memoria está completada en una primera versión y se encuentra en fase de revisión y mejora.
    \item Es necesario refinar contenido, especialmente a nivel de redacción y estructura.
\end{itemize}

\subsubsection*{2. Revisión del resumen e introducción}

\textbf{Observaciones:}
\begin{itemize}
    \item El resumen está demasiado centrado en el producto; debe incluir también información del proyecto (contexto académico, metodología, tecnologías, etc.).
    \item Se recomienda mantener una extensión breve (aprox. 3/4 de página) e incluir versión en inglés (abstract).
    \item La introducción puede repetir parcialmente el resumen, pero debe reforzar la motivación con datos económicos del dominio y añadir contexto que justifique la relevancia del proyecto.
\end{itemize}

\subsubsection*{3. Análisis y requisitos}

\textbf{Observaciones:}
\begin{itemize}
    \item Falta un apartado específico de actores del sistema (tipos de usuarios y diagrama de actores).
    \item Se deben añadir mockups como parte del análisis.
    \item Los requisitos funcionales deben reestructurarse en formato estándar: ``Como [usuario], quiero [funcionalidad], para [objetivo]''.
\end{itemize}

\subsubsection*{4. Objetivos del proyecto}
Se detecta mezcla de objetivos del sistema y objetivos del proyecto (aprendizaje, gestión, competencias). Se recomienda separarlos claramente y usarlos posteriormente en conclusiones para justificar resultados.

\subsubsection*{5. Metodología}
La sección actual mezcla distintos niveles. Se recomienda dividir en subsecciones: metodología de gestión del proyecto, metodología de desarrollo (sprints) y metodología técnica (gestión de código, commits, repositorio, etc.).

\subsubsection*{6. Planificación y fases}
\begin{itemize}
    \item Añadir una fase 0 (definición del alcance).
    \item Renombrar última fase a ``Despliegue y cierre del proyecto''.
    \item Incluir diagramas (Gantt u otros) y aportar datos cuantitativos (horas totales).
\end{itemize}

\subsubsection*{7. Presupuesto}
El coste actual es demasiado bajo y debe recalcularse utilizando datos oficiales, incluyendo: coste base por hora, costes sociales ($\sim$30\%) y costes indirectos (infraestructura). El presupuesto final debe incluir margen de contingencia (10--15\%) y beneficio industrial ($\sim$20\%).

\subsubsection*{8. Diseño del sistema}
\begin{itemize}
    \item Falta un diagrama conceptual independiente del tecnológico.
    \item Se debe crear un diagrama conceptual (sin detalles técnicos) junto al diagrama de diseño ya existente.
    \item Añadir sección de patrones de diseño y justificación de tecnologías utilizadas.
\end{itemize}

\subsubsection*{9. Diagramas y modelado}
Se recomienda incluir un diagrama de clases conceptual y diagramas UML claros, manteniendo coherencia en la representación.

\subsubsection*{10. Despliegue y uso real}
\begin{itemize}
    \item Backend desplegado; frontend desplegado sin publicación en stores (no es obligatorio publicar).
    \item Se recomienda añadir datos de uso, análisis de accesos y automatizaciones (cron jobs, scraping, etc.).
\end{itemize}

\subsubsection*{11. Anexos}
Se identifican elementos pendientes: actas de reuniones, informe de horas, resultados de pruebas y uso por usuarios piloto. Se recomienda incluir feedback de usuarios y evidencias del sistema en uso.

\subsubsection*{12. Entrega final}
Elementos necesarios para la entrega:
\begin{itemize}
    \item Código fuente (organizado).
    \item Aplicación funcional.
    \item Vídeo demostrativo ($\sim$3 minutos).
    \item Documentación completa.
\end{itemize}

\subsection*{Decisiones tomadas}
\begin{itemize}
    \item Reestructurar la memoria para mejorar claridad y separación de conceptos.
    \item Incluir actores, mockups y diagramas conceptuales.
    \item Recalcular completamente el presupuesto.
    \item Añadir métricas, datos reales y anexos faltantes.
    \item Reforzar la introducción con contexto económico.
\end{itemize}

\subsection*{Tareas asignadas}

\textbf{Estudiante (Cynthia):}
\begin{itemize}
    \item Revisar y mejorar el resumen e introducción.
    \item Separar objetivos del sistema y del proyecto.
    \item Añadir actores, mockups y diagramas conceptuales.
    \item Reestructurar la metodología.
    \item Recalcular el presupuesto.
    \item Completar anexos (actas, horas, pruebas, feedback usuarios).
    \item Preparar vídeo demostrativo.
    \item Finalizar documentación y presentación.
\end{itemize}

\textbf{Tutor (José Antonio):}
\begin{itemize}
    \item Revisar el documento completo.
    \item Proporcionar feedback adicional.
    \item Validar mejoras realizadas.
\end{itemize}

\subsection*{Próxima reunión}
\begin{itemize}
    \item \textbf{Fecha estimada:} en torno al 13 de mayo de 2026.
    \item \textbf{Objetivos:} Revisión final de la memoria, validación de correcciones y preparación de la entrega definitiva.
\end{itemize}

\textbf{Elaborado por:} Cynthia Castaño Juan \hfill \textbf{Fecha:} 6 de mayo de 2026
